\section{From Axis Mappings to Attribute Programs}
\label{sec:method}

Section~\ref{sec:mathematical-synthesis} defines the complete constructor.
This section compiles that same constructor into a symbolic metadata
program, before a match supplies concrete ranks or axis lists. Axis
mappings locate groups; partitioning separates their roles. The generated
program retains any group computations that are not simplified during
compilation. Native execution and checked symbolic lowering consume
the same program; they do not implement separate canonical constructors.

\subsection{Symbolic Axis Mapping}
\label{sec:symbolic-axis-mapping}
Section~\ref{sec:mathematical-synthesis} chooses logical axes for each RHS
operation, but operation attributes require zero-based dimension positions.
Compilation must express those positions using match metadata, before a
match supplies concrete ranks and axis lists. Symbolic mappings track where
variable-length groups occur and compose their position expressions into
RHS attributes.

A mapping associates a group $\vec G$ with a position-list expression $q_G$
in tensor $t$:
\begin{equation}
 \gath(\operatorname{Ord}_t,q_G)=\vec G.
 \label{eq:mapping}
\end{equation}
Here $\gath$ selects metadata-list entries at the specified positions.
$\operatorname{Ord}_t$ and $q_G$ are program expressions. At a match,
source replay assigns logical-axis identities from input-port connections;
equal extents never identify axes. Intermediate orders are expressions
constructed by Section~\ref{sec:mathematical-synthesis}'s passes.
Replaying source declarations supplies mappings from captured axis fields
and derived lists, such as ordered complements for free axes. For example,
$\operatorname{Ord}_T=[i_0,i_1,j_0,j_1]$ consists of groups
$\vec I=[i_0,i_1]$ and $\vec J=[j_0,j_1]$, with position lists
$q_I=[0,1]$ and $q_J=[2,3]$.

\noindent\textbf{Composing mappings.}
A selection $h$ within $\vec G$ has positions $\gath(q_G,h)$ in tensor
$t$, because
\begin{equation}
 \gath(\gath(\operatorname{Ord}_t,q_G),h)=\gath(\operatorname{Ord}_t,\gath(q_G,h)).
 \label{eq:gather-composition}
\end{equation}
Both selectors must be valid. For example, selecting $j_1$ within $\vec J$
uses $h=[1]$, giving $\gath(q_J,h)=\gath([2,3],[1])=[3]$ in $T$.

\subsection{Symbolic Axis Partitioning}
\label{sec:symbolic-axis-partitioning}
Continuing the example in Section~\ref{sec:symbolic-axis-mapping},
Figure~\ref{fig:axis-partitioning} shows a consumer contracting positions $d_\ell=[3,0]$ in $T$, paired
with $d_r=[2,0]$ in another operand $Z$. The selected axes $[j_1,i_0]$
cross the producer groups, so each group must be split into selected and
remaining axes.

\paperfigure{1}{figures/Symbolic Axis Partitioning (Dot Reassociation).pdf}{Symbolic axis partitioning for unbatched dot reassociation. Splitting selected axes by producer origin preserves their contraction partners. The figure uses $c_L,c_R$ for the text's $d_\ell,d_r$, $L$ for the selected contraction group, and $F=[f]$ for $Z$'s free group.}{fig:axis-partitioning}{Three panels show the intermediate tensor T with axes i0, i1, j0, j1; selection of positions 3 and 0; partition into selected and remaining I and J subgroups; and preservation of partners at positions 2 and 0 of Z. A table distinguishes positions within producer groups from slots in the contraction list and corresponding partner positions.}

For each $\vec G\in\{\vec I,\vec J\}$, compute
\begin{equation}
 (u_G,v_G)=\mpo(d_\ell,q_G),\qquad z_G=\gath(d_r,u_G).
 \label{eq:symbolic-partition-selectors}
\end{equation}
Paired matching finds common entries in $d_\ell$'s order, satisfying
$\gath(d_\ell,u_G)=\gath(q_G,v_G)$. Thus $u_G$ indexes slots in the
consumer's contraction list, whereas $v_G$ indexes axes within $\vec G$.
The selected subgroup is $\gath(\vec G,v_G)$; the ordered complement of
$v_G$ selects the remainder. These subgroups are disjoint and cover $\vec G$,
but concatenating them need not restore its order.

The figure's middle panel gives $(u_I,v_I)=([1],[0])$ and
$(u_J,v_J)=([0],[1])$. Applying these consumer slots to $d_r$ gives
$z_I=[0]$ and $z_J=[2]$, preserving the pairings shown in the right panel.
These pairings determine the replacement contractions; group positions
$v_G$ cannot substitute for consumer slots $u_G$.

\subsection{Symbolic Attribute-Program Construction}
\label{sec:symbolic-program-construction}
Mappings and partitions represent the sets and ordered groups of
Section~\ref{sec:mathematical-synthesis}. The compiler appends expressions
for its three passes: canonical retained axes, required-order checks, and
positional attributes. For example, mandatory axes are represented by
\begin{equation}
 M_v=\left(\bigcup_{X\in U_v}\operatorname{set}(I_X)\right)
 \cap\left(\operatorname{set}(F)\cup
             \bigcup_{X\notin U_v}\operatorname{set}(I_X)\right).
 \label{eq:emitted-mandatory}
\end{equation}
The subtree input sets $U_v$ are known at compilation; the source frames
$I_X$ and $F$ are supplied by replay at a match. Intersections and
differences implement the canonical role partitions; ordered filters and
concatenation implement the chosen layouts. Position queries convert the
resulting groups to RHS attributes. The same group order is used on both
sides of each paired contraction.

\noindent\textbf{One symbolic program.}
The compiler translates Section~\ref{sec:mathematical-synthesis}'s construction
into one program containing feasibility checks and computations of RHS
attributes. The mappings and partitions above express these computations
without concrete ranks or axis lists. Shared axis-group computations are
represented once and reused across attributes. This program is the compiler's
metadata intermediate representation (IR).

At compile time, TEPL lowers this symbolic metadata program into expressions
over source attributes, using axis mappings, partitioning, and simplification. For example, the selected subgroup $\gath(\vec J,v_J)$ has
positions $\gath(q_J,v_J)$ in $T$. The mapping for $\vec J$ yields this position
expression. At matching time, match metadata
supplies the lists: $q_J=[2,3]$ and $v_J=[1]$ give the position list $[3]$.
Thus compilation derives the formula; matching evaluates it to obtain the
RHS attribute. These computations operate on axis metadata, not tensor values.
A fresh validator checks the lowered expressions, their definedness,
and the rewrite obligations under the written guard.

Program construction does not require every computation to simplify.
If bounded symbolic lowering fails, the shared program is retained, but
Rust/C++ code generation requires this lowering to succeed. Direct
interpretation of the retained program is described in
Section~\ref{sec:implementation-scope}.

\begin{theorem}[Symbolic compilation correspondence]
\label{thm:symbolic-compilation}
For every admitted semantic plan in Section~\ref{sec:formal-scope}, compilation
constructs a finite metadata program $(E_q,E_f)$ independently of concrete
ranks and the written guard. For every legal metadata instance $m$ in that
section's completeness domain,
\begin{align}
 \lbrack\!\lbrack E_q\rbrack\!\rbrack_m&=\operatorname{check}(m),
 \label{eq:emitted-check}\\
 \operatorname{check}(m)&\Rightarrow
 \lbrack\!\lbrack E_f\rbrack\!\rbrack_m=\operatorname{build}_{\tau}(m).
 \label{eq:emitted-build}
\end{align}
Here $\tau$ fixes unconstrained paired orders using the left operand's
order. The predicate is total on the domain; builders are defined when it
holds. Sufficient resources and representable metadata are assumed.
\end{theorem}
\begin{proof}[Proof sketch]
The compiler translates each step of Section~\ref{sec:mathematical-synthesis}
into metadata operations that compute the same axis sets, ordering checks,
and positional attributes. Appendix~\ref{app:symbolic-compilation-proof}
proves this correspondence and shows that program construction terminates
without concrete metadata, the written guard, or successful simplification.
\end{proof}
Together with Theorem~\ref{thm:complete-synthesis}, this establishes
completeness of the generated feasibility/builder program on the stated
domain. Guard certification and checked symbolic lowering have separate
acceptance obligations; neither is claimed complete.

Appendix~\ref{app:batched-reassociation-program} gives the complete generated
attribute program for guarded batched reassociation.

\subsection{Acceptance and Implementation Boundary}
\label{sec:compilation-scope}
Table~\ref{tab:guarantee-map} collects the guarantees and their domains.
\begin{table}[!htbp]
\caption{Guarantees and their scope.}
\label{tab:guarantee-map}
\small
\begin{tabularx}{\linewidth}{@{}lY@{}}\toprule
Result & Guarantee\\\midrule
Constructor (Theorem~\ref{thm:complete-synthesis}) & Complete under D1--D4, read-once components, exact arithmetic, and extents $\ge2$; all RHS axis fields unknown.\\
Compilation (Theorem~\ref{thm:symbolic-compilation}) & Shared metadata IR realizes the constructor for admitted plans; compilation does not discharge guards.\\
Certified acceptance (Theorem~\ref{thm:accepted-core}) & Sound native portable execution under a certified guard; extents $\ge0$. Guard certification is sufficient, not complete.\\
Validation (Theorem~\ref{thm:validator-soundness}) & Sound completed candidates, including supplied fields, within the calculus; extents $\ge0$.\\
Certified schemas (Lemma~\ref{lem:certified-extensions}) & Sound broadcast/reverse, product, negation, and unit-view schemas, including certified supplied fields; no completeness claim.\\
Default extensions & Broader reasoning outside the proved fragments has empirical support only.\\\bottomrule
\end{tabularx}
\end{table}
Program construction and acceptance are separate. Acceptance requires
certification under the written guard for the whole-axis constructor or
the sound schemas in Appendix~\ref{app:axis-extensions}.
Schemas may preserve supplied fields when their equality with generated
metadata is certified; other candidates use the independent validator.

\paragraph{A small guard checker.}
Guard certification uses sufficient, explicit checks rather than
asking the general normalizer to simplify the entire generated program.
First, it recognizes five source/target structures with constructive
feasibility arguments: sum movement across contraction; transpose folding into contraction with
one free axis on the folded side; unary chains of
sums and permutations with sufficient RHS operations; equal expanded
component trees; and batched reassociation under the prefix-routing guard
of Appendix~\ref{app:batched-reassociation-program}.
Second, when input ranks and every source axis list are fixed by capture
constraints or positive conjuncts of $C$, it evaluates $E_q$ on that axis
layout with every extent set to two. This is an exact abstraction of
\emph{axis topology}: replay's port classes, ordered interfaces, and every
node of $E_q$ depend only on ranks and axis lists. No tensor values are
sampled. Therefore one successful abstract evaluation proves $E_q$ for
all legal extent assignments having that layout. The checker extracts no
facts from disjunctions or unmet obligations and adds no guard.
Appendix~\ref{app:accepted-core-proof} specifies these checks and their
constructive arguments.

The certified backend checks identity dtype transport for views, numeric
dtype preservation for sums and negation, and equal numeric input dtypes
for contraction and pointwise addition/product. Target numeric requirements
must follow from source legality. Custom dtype behavior is conservatively rejected.

\begin{theorem}[End-to-end certified acceptance]
\label{thm:accepted-core}
Assuming the trusted contracts in Section~\ref{sec:implementation-scope}, if
a rule satisfies the certification checks described above, native execution of its successfully
loaded program computes a legal replacement preserving values and ordered
shape for every legal source match satisfying its original guard $C$.
This holds under declared algebraic semantics, for nonnegative extents,
sufficient resources, and representable metadata.
\end{theorem}
\begin{proof}
For canonical emission, acceptance requires template admission, independent
automatic RHS fields, the dtype checks above, and guard certification. The
guard-checker lemmas establish $C\Rightarrow E_q$ from source legality.
Canonical dtype contracts preserve a common input dtype; source numeric
legality discharges every target numeric requirement.
The builder induction in Theorem~\ref{thm:symbolic-compilation} uses membership,
ordering, and paired extent equality, none of which requires a lower
extent bound. Its value argument uses finite sums and products and remains
valid for empty and singleton domains. Thus a true predicate suffices for
defined attributes, legal operations, and semantic equality. Only the
\emph{necessity} direction of completeness requires the minimum-extent
assumption; no completeness claim is made here for degenerate extents.
For the additional schemas, Lemma~\ref{lem:certified-extensions} supplies
the corresponding constructor argument. Guard certification checks its
conditions and equality of any supplied fields with the generated fields.
Materialization and loading preserve the checked DAG, schedule, and guard;
native evaluation executes the same metadata primitives.
\end{proof}

Portable programs store the canonical or schema DAG and its source/target schedule.
The loader regenerates the DAG, checks exact equality, binds the schedule
to actual operation contracts and fields, and reruns the guard checker on
the serialized conditions. The native evaluator caches the DAG when loading
and evaluates source replay and its builders per match. A failed canonical
predicate is an invariant failure, not an additional applicability guard;
resource exhaustion is reported separately. The loader checks this metadata
certificate; declaration semantics and their agreement with an execution
backend remain trusted.

These checks do not decide arbitrary guard implications. An unresolved
certification proof returns \code{Unknown}, without concluding infeasibility.
The validator checks inferred or supplied expression candidates. The
ordinary ordered strategy symbolically lowers the same canonical IR into
attribute expressions; generated Rust/C++ still require that lowering to succeed. Native
portable execution retains the residual IR. Broader general construction and
sound schema extensions remain outside the complete-constructor theorem.
